% Generated by records_reorg.py assemble (records-reorg/); edit the data and re-run, not this file.
\subsection{The estimator and its binned limit}\label{recsec:negweight-estimator}

% Moved verbatim from the note's App. B (app_negweight.tex) in the compression (B1; move-manifest-je JE-13).
% [records reorg] former heading: \subsection{Background-subtraction detail moved from the note}
\paragraph{Background-subtraction detail moved from the note.}

% [records reorg] former heading: \paragraph{Binned limit: the per-bin derivation (from App. B, sec:negweight-binned).}
\paragraph{Binned limit: the per-bin derivation (from App. B, sec:negweight-binned).} A step-1 classifier restricted to functions piecewise-constant on reconstructed analysis bins $\{b\}$ gives the population-level binned limit of the OmniFold/Richardson--Lucy correspondence in Refs.~\cite{Andreassen:2019cjw,Andreassen:2021zzk}. Its output is a single $p$ per bin. With common support and unnormalized class weights, step~1 weighted cross-entropy separates into per-bin objectives
\begin{equation}
  \label{eq:negweight-Lb}
  L_b(p) = -W_1\,\ln p - W_0\,\ln(1-p),
\end{equation}
where $W_0=S(b)>0$ is the integral of simulated class-0 density $S$ in bin $b$, and $W_1$ is the signed class-1 weight. Injecting background simulation at negative weight while keeping data at $+1$ gives $W_1=D(b)-B(b)$. Here $D(b)$ and $B(b)$ are the bin masses of $D$ and $B$, equivalently $N_{\mathrm{data}}(b)$ and $N_{\mathrm{bkg}}(b)$ in \S\ref{sec:negweight-estimator}.
For $D(b)>B(b)$, objective~\eqref{eq:negweight-Lb} is strictly convex on $(0,1)$ with unique interior minimizer $p=W_1/(W_0+W_1)$, giving
\begin{equation}
  % \label{eq:negweight-rb} is defined in the note (App. B keeps the equation); this copy drops the line.
  r(b) = \frac{p}{1-p} = \frac{W_1}{W_0} = \frac{D(b)-B(b)}{S(b)}.
\end{equation}

% [records reorg] former heading: \paragraph{Binned limit: the divergence for D(b)<B(b) (from App. B).}
\paragraph{Binned limit: the divergence for D(b)<B(b) (from App. B).} The reduction fails where $D(b)<B(b)$. With $W_1<0$, $L_b(p)=|W_1|\,\ln p-W_0\,\ln(1-p)$ tends to $-\infty$ as $p\to0^+$ and is unbounded below, with no minimizer.

% [records reorg] former heading: \subsubsection{Estimator and binned limit}
\paragraph{Estimator and binned limit.}
Negative-weight injection is the continuous version of the purity weight recommended for nontrivial backgrounds in Ref.~\cite{Canelli:2025guide}, and the natural treatment for the larger backgrounds of the extended-fiducial analysis.
A bin-constant rescaling reproduces $D(x)-B(x)$ inside bin $b$ only when $B/D$ is constant there. Otherwise, $D(x)\,(1-B_{b}/D_{b})$ and $D(x)-B(x)$ share a bin integral but differ within the bin. The injected background and data occupy the same reconstructed fiducial window, giving identical support.
On the exact-split backend in the backend comparison of App.~\ref{app:negweight} (\S\ref{sec:negweight-estimator}), the clipped refinement removes the instability that raw signed weights cause there.
A finite regularized classifier's behavior depends on the estimator. Stay-Positive remains well behaved but can retain negative refined weights where the target density is genuinely negative. The zero clip used here is therefore an additional analysis prescription chosen to match the physical purity floor. Ref.~\cite{Nachman:2025staypositive} documents the same reweighting-loss divergence for negative targets; signed binary cross-entropy has an analogous pathology for sWeights~\cite{GlazierTyson:2024}.
The binned $D-B$ update, negative-weight classification of the measured side, and binned OmniFold/Richardson--Lucy correspondence already appear together in Ref.~\cite{Andreassen:2021zzk}. Equation~\eqref{eq:negweight-rb} states their immediate binwise corollary and its precise validity domain for this analysis. It is a consistency check, not a separate method.

% [records reorg] former heading: \subsubsection{Closed-form toy}
\paragraph{Closed-form toy.}
A one-dimensional toy uses signal simulation $\mathcal{N}(0,1)$, equally normalized data signal $\mathcal{N}(0.3,1.1)$, and background $\mathcal{N}(2,0.7)$ contributing \SI{16.7}{\percent} of pseudo-data. The known target $(D-B)/S=\phi(x;0.3,1.1)/\phi(x;0,1)$ is smooth and positive. A signed-weight gradient-boosted classifier recovers it with median \SI{\nwToyNeg}{\percent} relative error, matching binned purity's \SI{\nwToyPur}{\percent} and reproducing it bin-for-bin on the same problem. Omitting subtraction leaves \SI{\nwToyBias}{\percent} residual-background bias. Results are stable across classifier seeds ($\pm\SI{\nwToySeed}{\percent}$), with no numerical instability from negative weights. Moving background into the sparse signal tail and deliberately over-subtracting makes the learned reweight floor cleanly at zero where $B>D$, as anticipated.

\subsection{Comparison, propagation and the production chains}\label{recsec:negweight-comparison}

% [records reorg] former heading: \subsubsection{Comparison on the 2D result and propagation}
\paragraph{Comparison on the 2D result and propagation.}
The 187-universe systematic ensemble varied genuine background consistently within each universe instead of freezing it at nominal; a statistical bootstrap used matched seeds. Agreement is at the level predicted by the $(D-B)/S$ identity. Residuals of \SI{\nwSystResid}{\percent} and \SI{\nwStatResid}{\percent} are consistent with ``to within \SI{2}{\percent}'' in \S\ref{sec:method}. Those residuals measure a close aggregate size for the tested constructions. They do not establish that the difference is sampling noise, and they do not establish the absence of a method bias. Negative-weight injection reproduces the central cross section \emph{per bin} (\S\ref{sec:negweight-2d}) and the covariance \emph{in aggregate}.

% [records reorg] former heading: \subsubsection{Which background treatment each production chain uses}
\paragraph{Which background treatment each production chain uses.}
Current production writes the background treatment into the record of every endpoint unfold as it is produced, so the background treatment of each new component can be read from the component itself; this changes how the treatment is recorded, not the treatment itself.
The standard chain uses purity on two grounds. First, a covariance that mixes background treatments is inferior to one that uses either treatment consistently: a lateral-band block built with negative-weight subtraction would be inconsistent with central and vertical-band unfolds built with purity subtraction. Second, switching would have required rerunning every unfold of the standard chain (its central, vertical-band, detector and endpoint unfolds) and would have invalidated the flux-corrected covariance then just chosen for use; this consequence, rather than the computing cost, decided the choice.
The cells where the treatments' mean difference exceeds $1\,\sigma_{\mathrm{stat}}$ in the 2D closure test with backgrounds are the whole lowest-$\pPar$ row ($1.5$--\SI{2.0}{GeV/c}) and five cells along the high-$\pT$, low-$\pPar$ boundary. In them, the purity estimate's mean residual against the truth is $+0.78$ to $+5.89\,\sigma_{\mathrm{stat}}$, above $1\,\sigma_{\mathrm{total}}$ in one cell ($1.22$); the refined estimate's mean residual is within $\pm0.44\,\sigma_{\mathrm{stat}}$ in every cell.
The integrated difference between the treatments is \SI{-0.198}{\percent}. In the median cell, the pseudo-experiments' own spread is $0.46\,\sigma_{\mathrm{stat}}$ (purity) and $0.58\,\sigma_{\mathrm{stat}}$ (refined).
The extended-fiducial battery's central value, \SI{\fpsSigma}{cm^2/nucleon}, is recorded with its passed published-region check against the control.

% [records reorg] former heading: \subsubsection{Displacement argument and scope}
\paragraph{Displacement argument and scope.}

% [records reorg] former heading: \paragraph{Displacement across universes.}
\paragraph{Displacement across universes.} The displacement argument rests on a \emph{measured} input; it assumes no identity of continuous estimands. The constructions share bin masses in the binned limit (\S\ref{sec:negweight-binned}) but differ by within-bin shape terms (\S\ref{sec:negweight-estimator}). Switching background treatment produces a measured net per-universe displacement of order \SI{0.1}{\percent}, including residual and shape effects. Systematic covariance measures spread across universes about their own mean. Similar small shifts of each universe leave that spread nearly unchanged, explaining the $187$-universe ratio $\nwSystRatio$, far from the purity correction's ${\sim}\SI{3}{\percent}$. This argument connects a handful-of-universes comparison to the $187$-universe covariance; the full $187$-universe rebuild in \S\ref{sec:negweight-uq} tests it at 2D and passes.
Raw negative weights destabilize the exact-split backend (\S\ref{sec:negweight-estimator}), but the standard chain uses LightGBM, so neither treatment is exposed to that instability here.

\subsection{Refined subtraction in five dimensions}\label{recsec:negweight-refined}

% [records reorg] former heading: \subsubsection{Refined subtraction in five dimensions}
\paragraph{Refined subtraction in five dimensions.}
The five-dimensional pseudo-experiments with both background treatments, and the purity bias they show at high $W$, are reported with Fig.~\ref{fig:bkg-residual-note}. This subsection describes the refined subtraction used in those tests, a refinement of higher capacity, and the independence of the background treatment from the choice of covariance.
% Received from cov2d (move-manifest-cv-negweight-return.tsv, CV-N01..N19, sec_q3_negweight.tex @ f6a4fc372; the director's routing, 2026-11-19):
% the three method paragraphs return to App. B verbatim; the result and its figure stay in Part 5.

% [records reorg] former heading: \paragraph{Refined negative-weight subtraction.}
\paragraph{Refined negative-weight subtraction.} Refinement used the unfold's deterministic GBDT configuration, independent injection and subtraction halves of background simulation, and bootstrap replicas that individually resampled every pseudo-event and refit refinement. It preserves the signed total to $3\times10^{-6}$ and clips \SIrange{0.01}{0.05}{\percent} of events. Corner bias falls from \SI{1.66}{\percent} to \SI{0.24}{\percent}, but residuals in some joint cells fail the nominal closure criterion. The criterion has two parts, each applied to all 153 reported quantities. No quantity may have a mean residual beyond 3.6 standard errors of the mean ($|t|>3.6$), a fixed threshold at about the two-sided 5\% level for a normal reference, Bonferroni-corrected over the 153. Nor may any quantity have a mean residual beyond half its bootstrap standard deviation with $|t|>3$. The data difference's sign pattern matches simulated purity bias at an order of magnitude smaller size; the difference is method sensitivity, not a measured bias. With either treatment, the estimator does not follow generator-sized $\Eavail$ shape changes. This refinement study ran no coverage test. Reducing one bias does not itself qualify an uncertainty, so this result does not close the revisit of the background treatment of the \resFiveAxis{} (\S\ref{sec:negweight-footing}).
% Note register round n1, batch n1-F: the criterion stated (state/s5n/contract.json, stage_1_development, controls C3
% and C6; OUTCOME-20260925-s5n-stage1-development-fail.md sec. 2, C3 FAIL; nd-unfolding/s5n_analyze_dev.py:112-149);
% "predeclared" dropped, since the result is a failure.
% Search: s5n coverage; coverage test; coverage verdict; no validation seed was run (OUTCOME-20260925-s5n-stage1-development-fail.md:13) @ ddf16f1087d7
% Source: docs/orchestration/OUTCOME-20260925-s5n-stage1-development-fail.md; VL151, VL152.
% Heading date (2026-09-25) and numbers repeated in sec_validation.tex removed (lane 3 App. B #7):
% highest-$W$ bias falls from \SIrange{-3.5}{-4.1}{\percent} to \SIrange{0.11}{0.24}{\percent};
% residual of up to about \SI{1}{\percent}; on the data the two footings
% differ by a median of \SI{0.36}{\percent} and at most \SI{1.4}{\percent}.

% [records reorg] former heading: \paragraph{Refinement capacity.}
\paragraph{Refinement capacity.} Stay-Positive classifier underfitting causes the residual nominal bias of the refined subtraction. At the unfold's GBDT setting (100 trees, 8 leaves), the positive-weight probability is underfit, leaving up to \SI{4}{\percent} reco-level subtraction error per $(\Eavail,W)$ reco cell. With 400 trees and 31 leaves, error falls to \SI{1.7}{\percent}, paired truth-level bias falls by more than \SI{80}{\percent} in every affected cell, and nominal background-inclusive closure passes on 20 development and 40 fresh pseudo-experiments. Splitting background simulation into independent injection and subtraction halves is not the cause.
% Source: docs/orchestration/DIAGNOSIS-20260926-s5e-oi192-estimator.md §5;
% docs/orchestration/OUTCOME-20260926-s5e-oi192-diagnosis-and-candidate.md; VL153, VL154.
% Heading date (2026-09-26) removed. Moved to the pointer above: "The higher-capacity refinement
% is not adopted: the estimator it feeds failed a stability test on the data (\S\ref{sec:validation})."

% [records reorg] former heading: \paragraph{Background treatment and choice of covariance.}
\paragraph{Background treatment and choice of covariance.} The choice of background treatment is independent of the choice of covariance, and App.~\ref{app:negweight} selects no five-dimensional covariance for use. The lateral-band component from re-selected event loops is built, with five bands and ten endpoint unfolds whose distinct checksums are listed in the record of the five-axis covariance of the \resFiveAxis{}. The status of that five-axis covariance and the scope of its use are stated in \S\ref{sec:eavailw-adopted}. The choice of background treatment settles one of the conditions for using a component; every other condition still applies.
% Plain status of the use of the five-axis covariance as printed at sec:eavailw-adopted (sec_5d.tex): it does not meet
% its specification's acceptance criteria, and its use covers only this file and its projection.
% Note consolidation (DECISION-20261027 item 1, cost-list R05): the status and the use's scope are no longer restated here; the sentence points to their single home, sec:eavailw-adopted, where the four measurements also print. "Neither fact" now refers to that status and that scope.
